\documentclass[../DoAn.tex]{subfiles}
\begin{document}

\section{Đặt vấn đề}
\label{section:1.1}

Ô nhiễm không khí là một trong những mối đe dọa sức khỏe cộng đồng nghiêm trọng nhất trên toàn cầu. Theo báo cáo năm 2021 của Tổ chức Y tế Thế giới (WHO), ô nhiễm không khí gây ra khoảng bảy triệu ca tử vong sớm mỗi năm, trong đó bụi mịn PM2.5 là tác nhân chính do khả năng xâm nhập sâu vào hệ hô hấp và tuần hoàn \cite{who2021}. Tại khu vực Đông Nam Á, Việt Nam là quốc gia chịu ảnh hưởng nặng nề, xếp thứ 23 trên thế giới và thứ hai trong khu vực về mức độ ô nhiễm không khí theo báo cáo chất lượng không khí toàn cầu năm 2024 của IQAir, với nồng độ PM2.5 trung bình năm đạt 28,7 $\mu g/m^3$ — cao gấp gần sáu lần so với ngưỡng khuyến nghị 5 $\mu g/m^3$ của WHO \cite{iqair2024report}. Riêng Hà Nội ghi nhận nồng độ PM2.5 trung bình năm là 45,4 $\mu g/m^3$, gấp chín lần ngưỡng WHO, khiến thành phố này xếp thứ bảy trong số 121 thủ đô được khảo sát \cite{iqair2024report}.

Trong bối cảnh đó, nhu cầu giám sát chất lượng không khí ở mức vi mô — từng khu phố, trường học, bệnh viện — ngày càng trở nên cấp thiết. Hệ thống quan trắc quốc gia VN-AQI hiện chỉ có khoảng 35 trạm đo tự động trên cả nước \cite{btnmt2023}, không đủ phản ánh sự biến thiên chất lượng không khí giữa các khu vực đô thị liền kề. Morawska và cộng sự đã chỉ ra rằng mạng lưới cảm biến chi phí thấp, mật độ cao là hướng tiếp cận khả thi để bổ sung cho các trạm đo chuẩn, qua đó cung cấp bức tranh toàn diện hơn về chất lượng không khí mà người dân tiếp xúc hàng ngày \cite{morawska2018}.

Tuy nhiên, các hệ thống giám sát chất lượng không khí hiện có — cả thương mại lẫn cộng đồng — đều có những hạn chế nhất định. Phần lớn các giải pháp như PAM Air, IQAir, PurpleAir yêu cầu kết nối WiFi hoặc 4G tại mỗi điểm đo, hạn chế khả năng triển khai tại các khu vực không có hạ tầng mạng sẵn có. Chi phí thiết bị thương mại cũng là rào cản lớn khi triển khai diện rộng, với mức giá từ vài triệu đến hàng trăm triệu đồng mỗi trạm. Ngoài ra, hầu hết các hệ thống đều là mã nguồn đóng, không cho phép tùy biến logic nghiệp vụ hoặc tích hợp với hệ thống bên thứ ba. Những hạn chế này đặt ra bài toán cần giải quyết: xây dựng một hệ thống giám sát chất lượng không khí đa thông số, chi phí thấp, có khả năng truyền dữ liệu tầm xa mà không cần hạ tầng WiFi tại mỗi điểm đo, và công khai toàn bộ mã nguồn.

\section{Mục tiêu và phạm vi đề tài}
\label{section:1.2}

Hiện nay có bốn nhóm giải pháp giám sát chất lượng không khí tiêu biểu. PAM Air triển khai mạng lưới hơn 3.000 trạm đo tại Việt Nam nhưng sử dụng mã nguồn đóng và chỉ đo PM2.5 \cite{pamair2024}. IQAir cung cấp nền tảng quốc tế với khả năng dự báo AI nhưng chi phí thiết bị cao (khoảng 269 USD mỗi thiết bị) và phần mềm không phải mã nguồn mở \cite{iqair2024}. PurpleAir dựa trên mô hình cộng đồng và cung cấp API mở nhưng chỉ đo bụi mịn, yêu cầu WiFi cho mỗi trạm \cite{purpleair2024}. Hệ thống quan trắc quốc gia VN-AQI đạt tiêu chuẩn FEM nhưng chi phí mỗi trạm từ hàng trăm triệu đến hàng tỷ đồng, số lượng trạm rất hạn chế \cite{btnmt2023}. Bên cạnh đó, xu hướng sử dụng cảm biến chi phí thấp kết hợp giao thức truyền thông công suất thấp tầm xa đã được nhiều nghiên cứu quốc tế khẳng định là hướng đi triển vọng \cite{morawska2018, augustin2016}.

Phân tích tổng quan cho thấy không giải pháp nào đồng thời đáp ứng bốn tiêu chí: (i) đo đa thông số môi trường, (ii) truyền dữ liệu tầm xa mà không cần WiFi tại từng điểm đo, (iii) mã nguồn mở hoàn toàn, và (iv) chi phí phần cứng thấp. Đặc biệt, không hệ thống nào sử dụng giao thức \gls{LoRa} — giao thức đã được chứng minh có khả năng truyền thông đến 10--15 km trong điều kiện tầm nhìn thẳng với mức tiêu thụ năng lượng rất thấp \cite{augustin2016}.

Trên cơ sở đó, đồ án đặt mục tiêu xây dựng một hệ thống giám sát chất lượng không khí end-to-end với năm chức năng chính: (i) thu thập dữ liệu sáu thông số môi trường (PM2.5, PM10, CO\textsubscript{2}, \gls{TVOC}, nhiệt độ, độ ẩm) từ các cảm biến chi phí thấp, (ii) truyền dữ liệu tầm xa qua \gls{LoRa} mà không cần WiFi tại từng trạm đo, (iii) xử lý, lưu trữ và tính chỉ số \gls{AQI} theo chuẩn \gls{USEPA} \cite{usepa2018}, (iv) hiển thị dashboard web thời gian thực cho người dùng gồm bản đồ AQI, biểu đồ lịch sử và xếp hạng ô nhiễm, và (v) cung cấp bảng quản trị cho quản trị viên gồm quản lý thiết bị, giám sát cảnh báo, cấu hình hệ thống, quản lý tài khoản và xuất dữ liệu. Hệ thống được phát triển dưới dạng mã nguồn mở hoàn toàn, cho phép bên thứ ba tái sử dụng, tùy biến và mở rộng.

\section{Định hướng giải pháp}
\label{section:1.3}

Để đáp ứng các mục tiêu đặt ra ở phần \ref{section:1.2}, em lựa chọn kiến trúc bốn tầng cho hệ thống. Tầng cảm biến sử dụng vi điều khiển ESP32 kết hợp ba module cảm biến chi phí thấp (PMS7003 đo bụi mịn, CCS811 đo CO\textsubscript{2} và TVOC, AHT10 đo nhiệt độ và độ ẩm). Firmware Sensor Node sử dụng cơ chế Deep-sleep để tối ưu năng lượng; trong phạm vi đồ án, prototype dùng board phát triển ESP32 DevKit V1 nên tuổi thọ pin đạt khoảng 20,5 giờ — hạn chế chủ yếu đến từ các linh kiện phụ trên board (chip USB-UART, LDO) chứ không phải từ kiến trúc firmware; với thiết kế phần cứng tùy chỉnh (bare-chip ESP32), thời gian hoạt động có thể đạt nhiều tháng. Tầng trung chuyển sử dụng Gateway ESP32 nhận dữ liệu qua giao thức \gls{LoRa} point-to-point ở tần số 433 MHz và chuyển tiếp lên máy chủ qua WiFi. Tầng máy chủ sử dụng Node.js/Express kết hợp PostgreSQL mở rộng bằng TimescaleDB (cơ sở dữ liệu chuỗi thời gian) và PostGIS (truy vấn không gian), giao tiếp thời gian thực qua Socket.IO. Tầng hiển thị sử dụng React/Vite kết hợp Leaflet (bản đồ) và Apache ECharts (biểu đồ). Toàn bộ hệ thống được container hóa bằng Docker để đơn giản hóa triển khai.

Sensor Node đo sáu thông số môi trường, đóng gói thành gói tin nhị phân 18 byte và gửi qua \gls{LoRa} đến Gateway. Gateway lưu gói tin vào bộ đệm vòng, sau đó gửi lên Backend Server dưới dạng HTTP POST. Backend Server xử lý dữ liệu theo pipeline ba bước: lọc trùng lặp, lưu vào cơ sở dữ liệu và tính chỉ số \gls{AQI}, kiểm tra ngưỡng cảnh báo và phát sóng dữ liệu mới đến giao diện web qua Socket.IO. Frontend cập nhật dashboard theo thời gian thực mà không cần người dùng tải lại trang.

Đồ án có bốn đóng góp chính. Thứ nhất, hệ thống \gls{IoT} đa tầng end-to-end hoạt động ổn định trong môi trường production, tích hợp bốn miền kỹ thuật khác biệt (phần cứng, firmware, backend, frontend) với tổng quy mô khoảng 14.036 dòng mã. Thứ hai, giao thức \gls{LoRa} binary 18 byte giảm thời gian chiếm dụng kênh truyền khoảng 11 lần so với \gls{JSON}, kèm cơ chế chống trùng lặp dữ liệu khi nhiều Gateway thu nhận cùng một gói tin. Thứ ba, pipeline xử lý telemetry trên Backend đảm bảo không mất dữ liệu đo ngay cả khi module cảnh báo xảy ra lỗi, với toàn bộ 44 kịch bản kiểm thử đạt kết quả pass. Thứ tư, chi phí phần cứng mỗi Sensor Node ở mức khoảng 875.000 VNĐ — thấp hơn giải pháp IQAir khoảng 11 lần — và toàn bộ mã nguồn được công khai.

\section{Bố cục đồ án}
\label{section:1.4}

Phần còn lại của báo cáo đồ án tốt nghiệp này được tổ chức như sau.

Chương 2 tiến hành khảo sát bốn hệ thống giám sát chất lượng không khí tiêu biểu (PAM Air, IQAir, PurpleAir, VN-AQI), phân tích ưu nhược điểm và xác định khoảng trống mà các giải pháp hiện tại chưa giải quyết được. Trên cơ sở đó, chương trình bày phân tích yêu cầu chức năng thông qua biểu đồ use case cho hai tác nhân Người dùng và Quản trị viên, đặc tả chi tiết năm use case quan trọng nhất, và xác định năm yêu cầu phi chức năng của hệ thống.

Chương 3 trình bày nền tảng lý thuyết và các công nghệ được lựa chọn cho từng tầng trong kiến trúc hệ thống. Nội dung bao gồm kiến trúc tổng thể bốn tầng, nền tảng phần cứng \gls{IoT} (vi điều khiển ESP32 và các module cảm biến), giao thức truyền thông \gls{LoRa}, công nghệ phần mềm phía máy chủ (Node.js, TimescaleDB, Socket.IO) và phía giao diện (React, Leaflet, ECharts), lý thuyết tính toán chỉ số \gls{AQI} theo công thức nội suy tuyến tính của \gls{USEPA}, và hạ tầng triển khai Docker.

Chương 4 trình bày chi tiết quá trình thiết kế kiến trúc (biểu đồ gói UML, biểu đồ trình tự), thiết kế phần cứng (sơ đồ mạch Sensor Node và Gateway), thiết kế firmware (giao thức \gls{LoRa} binary, luồng hoạt động deep sleep), thiết kế module Backend (pipeline xử lý telemetry, cơ chế cảnh báo), thiết kế cơ sở dữ liệu (TimescaleDB hypertable, Continuous Aggregate), và thiết kế giao diện người dùng. Chương cũng trình bày kết quả triển khai, kiểm thử (44/44 kịch bản pass, bao gồm 19 test case đơn vị module tính \gls{AQI}), và kết quả thực nghiệm trên môi trường production.

Chương 5 phân tích bốn đóng góp kỹ thuật nổi bật của đồ án, bao gồm việc xây dựng hệ thống \gls{IoT} đa tầng end-to-end, thiết kế giao thức \gls{LoRa} binary và cơ chế chống trùng lặp dữ liệu, kiến trúc Backend phân lớp với pipeline xử lý telemetry đảm bảo tính bền vững của dữ liệu, và giải pháp chi phí thấp với mã nguồn mở hoàn toàn. Mỗi đóng góp được phân tích theo cấu trúc bài toán, giải pháp và kết quả đạt được.

Chương 6 tổng kết toàn bộ đồ án, đánh giá mức độ hoàn thành so với mục tiêu đề ra, ghi nhận các hạn chế còn tồn tại (thời gian pin dài hạn, hiệu chuẩn cảm biến, thiếu ứng dụng di động), và đề xuất bốn hướng phát triển trong tương lai: tối ưu phần cứng, tích hợp chuẩn LoRaWAN, bổ sung mô hình dự báo \gls{AQI} và thông báo đẩy, cùng triển khai thử nghiệm quy mô lớn hơn.

\end{document}